% !TEX root = ../analisi-dei-capitolati.tex
% Capitolo 4 - C1 - AURA

\section{C1 -- AURA}\label{sec:c1}

\begin{scheda}
  \voce{Proponente}{Vimar S.p.A.}
  \voce{Dominio}{domotica e servizi di intelligenza artificiale generativa}
\end{scheda}

\subsection{Obiettivo}
Vimar chiede un server \glo{MCP} che renda lo standard \textbf{KNX IoT 3rd
Party API}, con cui si comandano i suoi impianti domotici, utilizzabile dai
principali servizi commerciali di intelligenza artificiale (Claude,
ChatGPT, Copilot, Gemini). L'utente può così interrogare e comandare la
propria casa in linguaggio naturale con l'assistente che preferisce.

\subsection{Requisiti principali}
Sono obbligatori tre componenti: un'infrastruttura a container descritta
come codice con Docker Compose e compatibile con AWS; un server MCP con
autenticazione OAuth2, selezione dell'impianto, lettura e scrittura sui
dispositivi, più un client MCP per i test; il collegamento a un servizio
di intelligenza artificiale commerciale, che deve gestire 17 casi d'uso
base e 21 avanzati. I test devono coprire almeno il 75\% del codice e
l'80\% dei requisiti end-to-end. Sono opzionali il rilascio su AWS, il
client via VS Code, CLI o TUI e la copertura oltre il 90\%; storico,
scenari composti e trigger sono \emph{nice-to-have}. Vimar fornisce un kit
di impianto fisico per gruppo e consiglia Python con l'SDK FastMCP.

\subsection{Incontro con il proponente}
Alla sessione congiunta del 22 ottobre 2026 (VE~2026-10-22) Vimar ha
chiarito che il linguaggio è libero, che la soluzione è principalmente
back-end e basta che giri in locale con Docker, che per lo stato il
\emph{polling} è accettato, che il kit resta a casa di un membro e si
comanda via cloud, e che il PoC minimo copre due o tre casi base. Nei test
il servizio di intelligenza artificiale va simulato, perché un LLM reale
non è deterministico.

\subsection{Valutazione del gruppo}

\begin{fortedeboli}
  \fd{\textbf{Tema innovativo}: MCP collegato a un impianto reale.}
     {\textbf{Hardware}: il kit fisico condiviso è un punto singolo di
      guasto.}
  \fd{\textbf{Requisiti molto dettagliati} e proponente molto
      strutturato.}
     {\textbf{Dominio da studiare}: KNX IoT e documentazione per
      integratori.}
  \fd{\textbf{Materiale fornito}: kit, API ed eventuali risorse AWS.}
     {\textbf{Parte conversazionale delegata} a servizi esterni.}
\end{fortedeboli}

C1 è la prima alternativa del gruppo: solido e ben seguito, ma con la
dipendenza da un dispositivo fisico e un interesse leggermente inferiore
per un progetto in cui l'intelligenza artificiale è affidata a servizi
commerciali.
